\documentclass[a4paper]{article}

\usepackage{graphicx}
\usepackage{amsmath,amssymb}
\usepackage{minted}
\usepackage{multirow}
\usepackage{float}
\usepackage{todonotes}
\usepackage{forest}
\usepackage{xstring}
\usepackage{xcolor}
\usepackage[most]{tcolorbox}
\usepackage{xparse}
\usepackage{natbib}

\usetikzlibrary{graphs,graphdrawing}
\usegdlibrary{layered}

% for external graphviz
\input{/home/donkarlo/Dropbox/repo/nd_ai_project/src/nd_ai/knoweldge_representation/language/formal/computer/programming/tex/latex/code_clips/external_graphviz.tex}

% for tags
\input{/home/donkarlo/Dropbox/repo/nd_ai_project/src/nd_ai/knoweldge_representation/language/formal/computer/programming/tex/latex/parts/control_sequence/kinds/semtag/semtag.tex}

% for links
\input{/home/donkarlo/Dropbox/repo/nd_ai_project/src/nd_ai/knoweldge_representation/language/formal/computer/programming/tex/latex/code_clips/seehere.tex}

\title{Optimal control theory}
\date{}

\begin{document}
    \maketitle
    \clickurlfooter{https://en.wikipedia.org/wiki/Optimal_control} Optimal control theory is a branch of control theory that deals with finding a control for a dynamical system over a period of time such that an objective function is optimized
    
    \cite{diedrichsen-2009-the-coordination-of-movement-optimal-feedback-control-and-beyond} explains how optimal control theory, especially optimal feedback control, helps account for flexible and task-dependent movement coordination. It argues that movement is organized by distributing control across multiple effectors, such as muscles, joints, and limbs, in a way that fits task demands.
    The authors show that this framework can explain patterns of feedback correction, movement variability, and coordination across effectors. They also point out that hierarchical control and movement initiation still need further development for a fuller theory of voluntary motor control.
    \bibliographystyle{plainnat}
    \bibliography{/home/donkarlo/Dropbox/repo/nd_shared_research/bibliography/bibliography.bib}
\end{document}